\begin{figure*}[!p]
\centering
\begin{minipage}[b]{.42\textwidth}
\includegraphics[width=1.0\textwidth]{2025_03/res/full_david_hawthorn}
\centerline{\emph{David's Hawthorn.}}
\end{minipage}\qquad
\begin{minipage}[b]{.42\textwidth}
\includegraphics[width=1.0\textwidth]{2025_03/res/full_nick_privet}
\centerline{\emph{Nick's Privet.}}
\end{minipage}

\begin{minipage}[b]{.42\textwidth}
\includegraphics[width=1.0\textwidth]{2025_03/res/full_brendan_yew}
\centerline{\emph{Brendan's Yew.}}
\end{minipage}\qquad
\begin{minipage}[b]{.42\textwidth}
\includegraphics[width=1.0\textwidth]{2025_03/res/full_tony_hawthorn}
\centerline{\emph{Tony's Hawthorn.}}
\end{minipage}

\begin{minipage}[b]{.42\textwidth}
\includegraphics[width=1.0\textwidth]{2025_03/res/full_terry_sweetgum}
\centerline{\emph{Terry's American Sweetgum.}}
\end{minipage}\qquad
\begin{minipage}[b]{.42\textwidth}
\includegraphics[width=1.0\textwidth]{2025_03/res/full_mark_hawthorn}
\centerline{\emph{Mark's Hawthorn.}}
\end{minipage}
\end{figure*}

\begin{figure*}[!p]
\centering
\begin{minipage}[b]{.42\textwidth}
\includegraphics[width=1.0\textwidth]{2025_03/res/full_steve_various}
\centerline{\emph{Steve's Elm, Hawthorn, and Cotoneaster composition.}}
\end{minipage}\qquad
\begin{minipage}[b]{.42\textwidth}
\includegraphics[width=1.0\textwidth]{2025_03/res/novice_raegan_larch}
\centerline{\emph{Raegan's Novice Larch.}}
\end{minipage}

\begin{minipage}[b]{.42\textwidth}
\includegraphics[width=1.0\textwidth]{2025_03/res/novice_stefano_cedar}
\centerline{\emph{Stefano's Novice Red Cedar.}}
\end{minipage}\qquad
\begin{minipage}[b]{.42\textwidth}
\includegraphics[width=1.0\textwidth]{2025_03/res/novice_edward_prunus}
\centerline{\emph{Edward's Novice Fuji Cherry.}}
\end{minipage}

\begin{minipage}[b]{.42\textwidth}
\includegraphics[width=1.0\textwidth]{2025_03/res/novice_edward_juniper}
\centerline{\emph{Edward's Novice Himalayan Juniper.}}
\end{minipage}\qquad
\begin{minipage}[b]{.42\textwidth}
\includegraphics[width=1.0\textwidth]{2025_03/res/novice_edward_hawthorn}
\centerline{\emph{Edward's Novice Paul's Scarlet Hawthorn.}}
\end{minipage}
\end{figure*}

\headline{{\textbf{\Large\sffamily The Top Chop:}}\\
          {\textbf{\large\sffamily The Tree of the Month Contest}}}
\label{2025-03-contest}
\noindent
The March meetup of our bonsai club was a thrilling showcase of imagination, skill, and the boundless creativity that defines our beloved art form.
With the theme \textit{``Yamadori, Urban or Wild,''} members brought forth an extraordinary array of trees, each telling a unique story of resilience and beauty.
The competition was fierce, with entries that resonated deeply with the theme, showcasing the artistry and dedication of our members.
The evening was a testament to the shared passion and skillfulness that make our club so special.

For those new to the contest, here's a quick reminder of the rules: each attendee can vote on every tree except their own.
Members award 1 to 3 marks for four aspects – trunk/branches, foliage, pot/container, and presentation (aligned with the month's theme) – with a maximum of 12 points per entry.
This structure ensures a fair and thoughtful evaluation of each tree's unique qualities.

\topic{Tree of the Month}
\noindent 
This month's entries in the \textit{Tree of the Month} category were nothing short of spectacular, with members pushing the boundaries of creativity and technique.
Here are the contenders, sorted by scored points:

\begin{enumerate}
    \item \textbf{David Fishwick}'s imposing hawthorn (\textit{Crataegus Monogyna}, 138 points) --- A breathtaking display of natural character, with rugged bark and dynamic movement that perfectly captured the essence of the theme.

    \item \textbf{Nick Lloyd}'s fissured privet (\textit{Ligustrum Vulgare}, 120 points) --- A striking composition, with deep vertical fissure and a sense of age that evoked the wild beauty of yamadori.

    \item \textbf{Brendan Gibbs}' towering yew (\textit{Taxus Baccata}, 110 points) --- A commanding presence, with a towering trunk and rich foliage that exuded timeless elegance.

    \item \textbf{Tony Ulatowski}'s slanted hawthorn (\textit{Crataegus Monogyna}, 110 points) --- A masterful slanting design, with delicate foliage and a sense of balance that drew admiration from all.

    \item \textbf{Terry Wilson}'s well\--balanced American sweetgum (Liquidambar Styraciflua, 107 points) – A harmonious composition, promising vibrant autumn colours and a sense of poise that delighted the eye.    

    \item \textbf{Mark Edwards}' contorted hawthorn (\textit{Crataegus Monogyna}, 97 points) --- A dramatic display of twisting branches and raw natural beauty, embodying the spirit of the theme.

    \item \textbf{Steve Palmer}'s inspiring composition with an elm, a hawthorn, and a cotoneaster (\textit{Ulmus Minor}, \textit{Crataegus Monogyna}, \textit{Cotoneaster Acuminatus}, 97 points) --- A creative group planting that told a story of inclusiveness and peace, earning praise for its originality.
\end{enumerate}

\noindent 
\textbf{David Fishwick} is the winner of the \textit{Tree of the Month for March 2025} with his massive hawthorn --- congratulations on this well\--deserved victory!

\topic{Novice Tree of the Month}
\noindent 
The \textit{Novice Tree of the Month} category is reserved for members who have been part of the club for less than two years, offering a platform to showcase their progress and celebrate their early achievements in the art of bonsai.
This month's entries were a joy to behold, sorted by scored points:

\begin{enumerate}
    \item \textbf{Raegan Gibbs}' exquisite European larch (\textit{Larix Decidua}, 97 points) --- A delicate and refined entry, with soft foliage and a sense of grace that hinted at great potential.

    \item \textbf{Stefano Bragaglia}'s rampant red cedar (\textit{Juniperus Virginiana}, 81 points) --- A dynamic composition, with vigorous growth and a sense of wild energy that captivated the audience.

    \item \textbf{Edward Mason}'s elegant Fuji cherry (\textit{Prunus Incisa}, 57 points) --- A charming tree, with dainty blossoms and a sense of elegance that showcased Edward's growing skill.
    
    \item \textbf{Edward Mason}'s charming Himalayan juniper (\textit{Juniperus Squamata}, 53 points) – A compact and well\--proportioned tree, with textured foliage and a sense of maturity beyond its years. 
    
    \item \textbf{Edward Mason}'s promising Paul's scarlet hawthorn (\textit{Crataegus Laevigata}, 33 points) --- A fruit that is still unripe but promises to yield succulent sweet fruits in terms of striking bold flowers as it continues to grow.
\end{enumerate}

\noindent
\textit{\textbf{Members are reminded that only one tree per meeting can accumulate points towards their tree of the month ranking.}}

\medskip
\textbf{Raegan Gibbs} is the winner of the \textit{Novice Tree of the Month for March 2025} --- well done on this fantastic achievement!

\begin{center}
    $\cdots$
\end{center}

\noindent
This month's competition showcased the limitless possibilities that bonsai embodies, reflecting enduring resilience, timeless elegance, and quiet strength.
Whether you're a seasoned practitioner or just starting your journey, there's always room to grow, learn, and draw inspiration from this art form.
Let's continue to challenge our creativity, share knowledge, and celebrate the accomplishments of our fellow enthusiasts.

Looking ahead, next month’s theme is \textit{``Juniper, Cypress, or Hemlock,''} celebrating these iconic trees and the diverse styles they can inspire---from formal upright to cascading, literati to forest plantings, and everything in between.
We're eager to witness the unique stories and personalities these trees will bring to our next gathering!

Well done to all participants, and thank you to everyone who voted.
Until next time, happy pruning!

\closearticle

\columnbreak
